% 【专题研讨】所属：第1章（由 chapters/revised/01-knowledge.tex \input 引入）
\minitopic{专题研讨：相信有规范吗}\index{认识规范}\index{信念伦理} 认识论不只描述人怎样形成信念，还评价怎样相信才合宜。但“应当相信$p$”至少有四种读法。它可能说证据支持$p$，说某种过程可靠，说主体在现有能力下无责，或说接受$p$有利于行动。四项常常同向，却能分离：一项证据很强的坏消息可能令人痛苦；一个无责主体可能因环境欺骗而形成不可靠信念；一个有利信念也可能没有证据。 规范的对象同样需要区分。我们似乎不能仅凭决定就立即相信“明天不会下雨”，却能决定是否查预报、是否只读支持自己的一方、是否记录反证。对信念缺少直接控制，不等于没有认识责任。规范可以评价当前态度与证据是否相称，也可以评价主体在注意、查询和维护上的间接控制。 如果认识义务要求“不相信任何证据不足的命题”，日常生活会不会过度苛刻？人无法核验每项常识，也必须依赖记忆和证言。较合理的义务应随问题重要性、核验成本、具体警报和主体角色变化。无具体风险时默认使用正常来源可以合理；发现来源异常仍拒绝复核，则产生新的责任问题。

\minitopic{真理是否是唯一认识目标}信念的正确性与真值密切相关：相信$p$时，主体把$p$当作真。如果认识规范完全不关心真，便难解释证据为何应引导信念。但从“真理是基本目标”不能推出“应尽可能拥有最多真信念”。记住海滩每粒沙的位置虽增加真信念，却挤占有限注意，也未必有理解价值。 一种方案把目标写成“相信真命题且不相信假命题”。这包含追求真与避免错误两方面，而两者可能冲突：只相信逻辑真理可以少错，却错过世界信息；轻易相信则获得一些真理，同时积累大量错误。具体规范需说明在何种任务中如何权衡信息增益与错误风险。 理解、智慧和认识自主也可能是不可还原的目标。理解要求把握关系，智慧要求判断重要性，自主要求管理依赖。它们通常促进真理，却不能只按真信念数量衡量。一个学生偶然背对全部答案，拥有许多真信念，却没有稳定能力；另一个提出错误假说后准确诊断失败，可能取得重要认识进步。 真理目标还受道德限制。未经同意获取他人隐私可能得到真信息，却不因此成为许可的探究。保护研究参与者、延迟公开危险信息和尊重保密，都不是把假当真，而是限制何种求真行动可以采取。评价时应保持认识理由与道德理由的种类，不用一个总分掩盖冲突。

\minitopic{合理信念是否唯一许可}给定同一证据，是否只有一个理性态度？严格唯一论认为证据决定一个信念或精确置信度；许可主义允许不同主体在相同证据下采取多个合理态度。争论触及背景信念、风险态度和认识目标能否合理不同。 若两位气候研究者掌握同一数据，却使用都受到学界支持的模型，得出60\%与75\%预测，差异可能来自模型不确定性，而非一人不理性。若允许范围过宽，许可主义会让证据失去约束；若要求唯一数字，理论又需说明有限数据如何决定全部先验和模型。 “同一证据”本身并不容易满足。两人可能看过同一文档，却在理解、记忆和背景专长上不同；也可能拥有相同命题材料，但某证据对其中一人更显著。思想实验可以规定完全相同，现实分析则应先排除隐藏差异。 分歧出现后，双方又获得关于对方判断的高阶证据。即使原先各自被许可，知道同侪选择不同态度也可能要求更新。许可主义不等于每个人永远坚持初始立场，它仍需说明反思与社会校准。

\minitopic{认识责任的四种失败}第一种是证据忽视：主体接触到明确反证，却因不喜欢结论而不纳入判断。第二种是选择性搜寻：只查询会确认原信念的来源。第三种是来源失管：知道渠道有系统错误仍继续无保留依赖。第四种是过度断言：把暂时、局部或概率证据说成确定普遍事实。 四种失败的心理原因不同。证据忽视可能是动机偏差，选择性搜寻可能由平台推荐放大，来源失管可能来自时间压力，过度断言可能受奖励机制驱动。把它们都归于“缺乏批判性思考”不能指导改进。对应措施分别是反方清单、预先搜索协议、来源审计和限定语言。 责任也有能力条件。儿童、信息被封锁者和专业外行不能承担与专家相同义务；但能力有限不意味着任何信念都合宜。主体可承担与角色相称的最低责任，例如承认不知道、求助可信来源和避免高风险断言。 制度会改变个人责任。若医生每次复核都会受惩罚，单纯要求“更认真”忽略激励；若平台隐藏来源，用户很难评估独立性。认识责任既可评价个人选择，也可评价谁设计了信息环境、谁控制时间和复核权限。

\minitopic{自然化解释与规范评价}自然化认识论研究实际认知机制、环境和错误分布。实验可发现记忆受提问影响、视觉可靠性随光线变化、群体判断受发言顺序影响。这些事实对可行规范至关重要：要求人执行心理上不可能的更新，或假定不存在的注意资源，不是好规范。 经验事实却不单独推出规范。人普遍确认偏误，不等于应该偏误；某策略在历史环境有效，也不证明在新平台合理。规范理论仍需目标，如真理、理解或可靠行动，并说明经验机制如何服务目标。蒯因的自然化方案推动从先验奠基转向研究感觉输入与理论输出，但规范维度是否被保留一直是争点\parencite{quine1969}。 一种合作模型分三步。哲学先澄清评价对象和目标；经验研究测量实际机制和环境；规范设计再比较可行干预。三步互相修正：实验证明案例直觉不稳定，概念分析需要调整证据权重；规范目标变化，又会改变实验应测什么。 自然主义还提醒理论的环境敏感性。一个启发式在信息稀缺时节约资源，在操纵平台上却被利用；面对熟人可靠的信任习惯，在匿名大规模传播中可能失败。不能给过程一个脱离环境的永久可靠标签。

\minitopic{比较视角：学、问与可教性}《论语》常把知与不知、学、问和改过放在关系实践中\parencite{analects2003}。这提供一种认识主体图景：好主体不仅拥有正确答案，还能够承认边界、向他人学习并修正。它与当代认识责任相通，却不等同一条关于证据充分性的形式原则。 “可教”具有双重面向。主体要有接受理由和纠正的开放，教师与制度也要值得信任并允许提问。把顺从权威当作可教，会牺牲自主；把拒绝一切依赖当作独立，又使学习不可能。关系质量是认识能力的组成条件。 古典修养语言还会把认识、品格与行动紧密连接。当代理论可以区分信念确证与人格德性，避免把一个正确但性格有缺陷的人说成完全无知识；修养论则反问，只评价孤立信念是否忽略了稳定偏差和长期学习。两种尺度可并存，但不可在论证中无提示切换。

\minitopic{综合研讨：校园健康信息争议}校园出现传染病传闻。匿名帖称已有多人确诊，校方公告说“尚无经实验室确认病例”，学生会要求立即停课。三个句子可能同时真实：有人出现症状，实验室尚未确认，预防行动仍有理由。把它们写成二选一真假争议，会混淆命题和行动门槛。 先建立证据表：匿名帖有何直接接触，校方“尚无确认”是否被误读为“没有病例”，检测是否延迟，学生会依据什么损失评估。再区分认识态度：可认为传播风险值得重视，对确诊人数悬置，对短期减少聚集采取预防。 责任评价也要分配。发帖者若编造数字负有真实性责任；校方若用技术表述掩盖已知症状负有沟通责任；学生转发时应保留不确定性；平台应把后续更正连接原帖。最终是否确有病例，不能回溯改变早期根据是否合宜。 研讨可由四组分别主张证据主义、可靠主义、德性认识论和社会制度方案。每组必须回答：目标态度是什么，决定因素在哪里，什么信息会击败它，个人与机构各负何责。最后不投票选“正确主义”，而比较哪一方案解释了哪些维度。

\begin{tcolorbox}[breakable,colback=white,colframe=blue!30!black,title={专题研讨任务},fonttitle=\bfseries]
围绕一项现实争论写四页研究备忘录：第一页区分真值、确证、责任和行动；第二页重构两种规范理论；第三页给出一个最小对比和一项经验研究设计；第四页说明个人与制度改进，并列出会改变你结论的证据。
\end{tcolorbox}
